\QHeading{第1题　在$\square$里填上合适的数或字母。}
\par\noindent


\Equation{125\times 32=(125\times\textcolor{answerred}{8})\times 4}
\par

\Equation{(a-b)\cdot c=\textcolor{answerred}{a}\cdot\textcolor{answerred}{c}-\textcolor{answerred}{b}\cdot\textcolor{answerred}{c}}
\Equation{98+63+37=98+(\textcolor{answerred}{63}+\textcolor{answerred}{37})}
\Equation{(15+\textcolor{answerred}{12})\times 4=\textcolor{answerred}{15}\times 4+12\times\textcolor{answerred}{4}}
\Equation{90\div 2\div 15=90\div (\textcolor{answerred}{2}\times 15)}
\Equation{469-142-58=469-(\textcolor{answerred}{142}+\textcolor{answerred}{58})}
\par\noindent

\QHeading{第2题　用简便方法计算下面各题。}

\CalcBlock[方法一：简便计算]{1}{\begin{aligned}[t]&\phantom{={}}\color{black}300-138-162\\&\color{answerred}=300-(138+162)\\&\color{answerred}=300-300\\&\color{answerred}=0\end{aligned}}
\par
\MethodBlock{方法二：先处理括号或乘除，再提取相同因数}{用乘法分配律反过来计算。}{\begin{aligned}[t]&\phantom{={}}300-138-162\\
&=162-162\\
&=162\times (1-1)\\
&=162\times 0\\
&=0\end{aligned}}


\CalcBlock[方法一：简便计算]{2}{\begin{aligned}[t]&\phantom{={}}\color{black}367+91+33\\&\color{answerred}=(367+33)+91\\&\color{answerred}=400+91\\&\color{answerred}=491\end{aligned}}
\MethodBlock{方法二：按运算顺序直接计算（对照）}{本法只作结果核对；本题要求简便计算，前面的简便解法仍保留。}{\begin{aligned}[t]&\phantom{={}}367+91+33\\
&=458+33\\
&=491\end{aligned}}

\CalcBlock[方法一：简便计算]{3}{\begin{aligned}[t]&\phantom{={}}\color{black}434+232+168\\&\color{answerred}=434+(232+168)\\&\color{answerred}=434+400\\&\color{answerred}=834\end{aligned}}
\MethodBlock{方法二：先处理括号或乘除，再凑整再减去多加的数}{}{\begin{aligned}[t]&\phantom{={}}434+232+168\\
&=666+168\\
&=666+170-2\\
&=836-2\\
&=834\end{aligned}}

\par\noindent

\QHeading{第3题　下面每组算式中，得数相等的在（　）里把得数写出来。}

\par\noindent\makebox[\linewidth][c]{
\begin{tikzpicture}[x=1bp,y=-1bp,color=black,draw=black,text=black]
\begin{scope}\clip (67.0,393.33) rectangle (522.0,559.33);
\hdhead{37}{运算律——复习}
\hdsection{1}{151}{在$\square$里填上合适的数或字母。}
\q{106}{205}{390}{12.5}{$125\times 32=(125\times $}
\q{331}{205}{390}{12.5}{$)\times 4$}
\q{469}{205}{390}{12.5}{$(a-b)\cdot c=$}
\q{647}{205}{390}{12.5}{$\cdot$}
\q{710}{205}{390}{12.5}{$-$}
\q{773}{205}{390}{12.5}{$\cdot$}
\q{106}{260}{390}{12.5}{$98+63+37=98+($}
\q{358}{260}{390}{12.5}{$+$}
\q{418}{260}{390}{12.5}{$)$}
\q{469}{260}{390}{12.5}{$(15+$}
\q{568}{260}{390}{12.5}{$)\times 4=$}
\q{685}{260}{390}{12.5}{$\times 4+12\times $}
\q{106}{315}{390}{12.5}{$90\div 2\div 15=90\div ($}
\q{345}{315}{390}{12.5}{$\times 15)$}
\q{469}{313}{390}{12.5}{$469-142-58=469-($}
\q{752}{313}{390}{12.5}{$+$}
\q{816}{313}{390}{12.5}{$)$}
\begin{scope}[scale=.6666667,draw=black,text=black]\draw[line width=.5pt] (291,197) rectangle (326,231);\end{scope}
\begin{scope}[scale=.6666667,draw=black,text=black]\draw[line width=.5pt] (606,195) rectangle (641,229);\end{scope}
\begin{scope}[scale=.6666667,draw=black,text=black]\draw[line width=.5pt] (669,195) rectangle (704,229);\end{scope}
\begin{scope}[scale=.6666667,draw=black,text=black]\draw[line width=.5pt] (732,195) rectangle (767,229);\end{scope}
\begin{scope}[scale=.6666667,draw=black,text=black]\draw[line width=.5pt] (795,195) rectangle (830,229);\end{scope}
\begin{scope}[scale=.6666667,draw=black,text=black]\draw[line width=.5pt] (315,252) rectangle (350,286);\end{scope}
\begin{scope}[scale=.6666667,draw=black,text=black]\draw[line width=.5pt] (380,252) rectangle (415,286);\end{scope}
\begin{scope}[scale=.6666667,draw=black,text=black]\draw[line width=.5pt] (529,250) rectangle (564,284);\end{scope}
\begin{scope}[scale=.6666667,draw=black,text=black]\draw[line width=.5pt] (642,250) rectangle (677,284);\end{scope}
\begin{scope}[scale=.6666667,draw=black,text=black]\draw[line width=.5pt] (789,250) rectangle (824,284);\end{scope}
\begin{scope}[scale=.6666667,draw=black,text=black]\draw[line width=.5pt] (304,307) rectangle (339,341);\end{scope}
\begin{scope}[scale=.6666667,draw=black,text=black]\draw[line width=.5pt] (711,306) rectangle (746,340);\end{scope}
\begin{scope}[scale=.6666667,draw=black,text=black]\draw[line width=.5pt] (776,306) rectangle (811,340);\end{scope}
\hdsection{2}{376}{用简便方法计算下面各题。}
\q{105}{420}{260}{13}{$300-138-162$}
\q{370}{420}{260}{13}{$367+91+33$}
\q{633}{420}{260}{13}{$434+232+168$}
\hdsection{3}{557}{下面每组算式中，得数相等的在（　）里把得数写出来。}
\q{105}{592}{260}{12.5}{$\left\{\begin{aligned}&810\div 9\div 3\\&810\div (9\times 3)\end{aligned}\right.$}
\q{364}{592}{260}{12.5}{$\left\{\begin{aligned}&43\times 101-43\\&43\times (101-1)\end{aligned}\right.$}
\q{618}{592}{260}{12.5}{$\left\{\begin{aligned}&900\div (45\div 5)\\&900\div 45\times 5\end{aligned}\right.$}
\q{105}{721}{260}{12.5}{$\left\{\begin{aligned}&600\div (12\times 5)\\&600\div 12\div 5\end{aligned}\right.$}
\q{364}{721}{260}{12.5}{$\left\{\begin{aligned}&125\times (40+8)\\&125\times 40+125\times 8\end{aligned}\right.$}
\q{618}{721}{260}{12.5}{$\left\{\begin{aligned}&560-(360+173)\\&560-360-173\end{aligned}\right.$}
\q{154}{685}{28}{13}{（}
\q{214}{685}{25}{13}{）}
\q{410}{685}{28}{13}{（}
\q{470}{685}{25}{13}{）}
\q{663}{685}{28}{13}{（}
\q{723}{685}{25}{13}{）}
\q{154}{815}{28}{13}{（}
\q{214}{815}{25}{13}{）}
\q{410}{815}{28}{13}{（}
\q{470}{815}{25}{13}{）}
\q{663}{815}{28}{13}{（}
\q{723}{815}{25}{13}{）}
\hdsection{4}{873}{解决问题。}
\q{107}{919}{740}{12}{王老师去采购乐器，买了5把小号和5支长笛。已知每把小号520元，}
\q{107}{961}{740}{12}{每支长笛480元。}
\q{107}{1006}{352}{12}{① 王老师带了9500元，还剩下}
\q{107}{1048}{352}{12}{　多少钱?}
\q{504}{1006}{334}{12}{② 买小号比买长笛多花多少钱?}

\begin{scope}[color=answerred,draw=answerred,text=answerred]% 高清专用：仅调整填空/数轴标签位置，计算内容沿用对应红笔片段。
\ans{202.00}{135.33}{33.33}{10}{$\vphantom{0}8$}
\ans{216.00}{172.67}{33.33}{10}{$\vphantom{0}63$}
\ans{259.33}{172.67}{33.33}{10}{$\vphantom{0}37$}
\ans{208.67}{209.33}{33.33}{10}{$\vphantom{0}2$}
\ans{410.00}{135.33}{33.33}{10}{$\vphantom{0}a$}
\ans{452.00}{135.33}{33.33}{10}{$\vphantom{0}c$}
\ans{493.33}{135.33}{33.33}{10}{$\vphantom{0}b$}
\ans{535.33}{135.33}{33.33}{10}{$\vphantom{0}c$}
\ans{358.67}{172.00}{33.33}{10}{$\vphantom{0}12$}
\ans{433.33}{172.00}{33.33}{10}{$\vphantom{0}15$}
\ans{532.00}{172.00}{33.33}{10}{$\vphantom{0}4$}
\ans{478.67}{208.00}{33.33}{10}{$\vphantom{0}142$}
\ans{522.00}{208.00}{33.33}{10}{$\vphantom{0}58$}
\ans{70.67}{304.67}{166.67}{9.5}{$\vphantom{0}\begin{aligned}[t]&=300-(138+162)\\&=300-300\\&=0\end{aligned}$}
\ans{246.67}{304.67}{166.67}{9.5}{$\vphantom{0}\begin{aligned}[t]&=(367+33)+91\\&=400+91\\&=491\end{aligned}$}
\ans{421.33}{304.67}{166.67}{9.5}{$\vphantom{0}\begin{aligned}[t]&=434+(232+168)\\&=434+400\\&=834\end{aligned}$}
\ans{116.00}{459.33}{33.33}{10}{$\vphantom{0}30$}
\ans{291.00}{459.33}{33.33}{10}{$\vphantom{0}4300$}
\ans{454.67}{459.33}{33.33}{10}{$\vphantom{0}100$}
\ans{115.33}{545.33}{33.33}{10}{$\vphantom{0}10$}
\ans{291.00}{545.33}{33.33}{10}{$\vphantom{0}6000$}
\ans{456.67}{545.33}{33.33}{10}{$\vphantom{0}27$}
\ans{72.00}{717.33}{243.33}{9.5}{$\setlength{\jot}{-2pt}\begin{aligned}[t]&\phantom{={}}9500-(520+480)\times5\\&=9500-1000\times5\\&=9500-5000\\&=4500\text{（元）}\end{aligned}$}
\ans{335.33}{694.00}{223.33}{9.5}{$\setlength{\jot}{-2pt}\begin{aligned}[t]&\phantom{={}}(520-480)\times5\\&=40\times5\\&=200\text{（元）}\end{aligned}$}
\ans{335.33}{740.00}{223.33}{9.5}{答：买小号比买长笛多花200元。}

\ans{72}{777}{243.33}{9.5}{答：王老师还剩4500元。}
\end{scope}
\end{scope}\pgfresetboundingbox\path[use as bounding box] (67.0,393.33) rectangle (522.0,559.33);
\end{tikzpicture}}\par



\par




\par\noindent

\Answer{
\begingroup\color{answerred}
\QHeading{计算过程：3. 判断每组是否相等并写得数}

每组的两个算式分别计算，再比较得数。

\noindent
（1）
\Formula{\begin{gathered}\begin{aligned}[t]&\phantom{={}}810\div9\div3\\&=90\div3\\&=30\end{aligned}\qquad\begin{aligned}[t]&\phantom{={}}810\div(9\times3)\\&=810\div27\\&=30\end{aligned}\end{gathered}}
两个得数相等，在括号里填30。
\par\medskip

\noindent
（2）
\Formula{\begin{gathered}\begin{aligned}[t]&\phantom{={}}43\times101-43\\&=43\times(101-1)\\&=43\times100\\&=4300\end{aligned}\qquad\begin{aligned}[t]&\phantom{={}}43\times(101-1)\\&=43\times100\\&=4300\end{aligned}\end{gathered}}
两个得数相等，在括号里填4300。
\par\medskip

\noindent
（3）
\Formula{\begin{gathered}\begin{aligned}[t]&\phantom{={}}900\div(45\div5)\\&=900\div9\\&=100\end{aligned}\qquad\begin{aligned}[t]&\phantom{={}}900\div45\times5\\&=20\times5\\&=100\end{aligned}\end{gathered}}
两个得数相等，在括号里填100。
\par\medskip

\noindent
（4）
\Formula{\begin{gathered}\begin{aligned}[t]&\phantom{={}}600\div(12\times5)\\&=600\div60\\&=10\end{aligned}\qquad\begin{aligned}[t]&\phantom{={}}600\div12\div5\\&=50\div5\\&=10\end{aligned}\end{gathered}}
两个得数相等，在括号里填10。
\par\medskip

\noindent
（5）
\Formula{\begin{gathered}\begin{aligned}[t]&\phantom{={}}125\times(40+8)\\&=125\times40+125\times8\\&=5000+1000\\&=6000\end{aligned}\qquad\begin{aligned}[t]&\phantom{={}}125\times40+125\times8\\&=5000+1000\\&=6000\end{aligned}\end{gathered}}
两个得数相等，在括号里填6000。
\par\medskip

\noindent
（6）
\Formula{\begin{gathered}\begin{aligned}[t]&\phantom{={}}560-(360+173)\\&=560-533\\&=27\end{aligned}\qquad\begin{aligned}[t]&\phantom{={}}560-360-173\\&=200-173\\&=27\end{aligned}\end{gathered}}
两个得数相等，在括号里填27。
\par\medskip


\endgroup

}

\QHeading{第4题　解决问题。}

\PlainItem{王老师去采购乐器，买了5把小号和5支长笛。已知每把小号520元，每支长笛480元。}
\par

\Problem{① 王老师带了9500元，还剩下多少钱?}

\MethodTitle{方法一}
\Answer{\Formula{\begin{aligned}[t]&\phantom{={}}9500-(520+480)\times5\\&=9500-1000\times5\\&=9500-5000\\&=4500\text{（元）}\end{aligned}}}
\Answer{答：王老师还剩4500元。}

\MethodTitle{方法二：分别扣去小号和长笛的钱}
\Answer{先扣去5把小号的钱，再扣去5支长笛的钱。}
\Answer{\Formula{\begin{aligned}[t]&\phantom{={}}9500-520\times 5-480\times 5\\
&=9500-2600-2400\\
&=6900-2400\\
&=4500\text{（元）}\end{aligned}}}
\Answer{答：王老师还剩4500元。}
\Problem{② 买小号比买长笛多花多少钱?}

\MethodTitle{方法一}
\Answer{\Formula{\begin{aligned}[t]&\phantom{={}}(520-480)\times5\\&=40\times5\\&=200\text{（元）}\end{aligned}}}
\Answer{答：买小号比买长笛多花200元。}

\MethodTitle{方法二：分别求两种乐器的总价再求差}
\Answer{两种乐器各买5件，先算各自总价。}
\Answer{\Formula{\begin{aligned}[t]&\phantom{={}}520\times 5-480\times 5\\
&=2600-2400\\
&=200\text{（元）}\end{aligned}}}
\Answer{答：买小号比买长笛多花200元。}

